% !TeX spellcheck = pl_PL
\documentclass[a4paper,twoside]{article}
\usepackage{polski}
\usepackage[utf8]{inputenc}
\usepackage[pdftex]{graphicx}
\usepackage{amsmath}

\usepackage[unicode, bookmarks=true]{hyperref} %do zakładek
\usepackage{tabto} % do tabulacji
\NumTabs{11} % globalne ustawienie wielkosci tabulacji
\usepackage{array}
\usepackage{multirow}
\usepackage{array}
\usepackage{dcolumn}
\usepackage{bigstrut}
\usepackage{color}
\usepackage[usenames,dvipsnames]{xcolor}
\usepackage{pdfpages}
\usepackage{sidecap}
\usepackage{wrapfig}
\usepackage{float}	%for figure& table placement in text
\usepackage{textgreek}


\setlength{\textheight}{24cm}
\setlength{\textwidth}{15.92cm}
\setlength{\footskip}{10mm}
\setlength{\oddsidemargin}{0mm}
\setlength{\evensidemargin}{0mm}
\setlength{\topmargin}{0mm}
\setlength{\headsep}{5mm}
\newcommand{\HRule}{\rule{\linewidth}{0.5mm}} 

\newcolumntype{M}[1]{>{\centering\arraybackslash}m{#1}}
\newcolumntype{N}{@{}m{0pt}@{}}

\graphicspath{ {./img/} }

% === Reset inkrementacji sekcji przy nowym parcie === %
\usepackage{titlesec}

\makeatletter
\@addtoreset{section}{part}
\makeatother
\titleformat{\part}[display]
{\normalfont\LARGE\bfseries\left}{}{0pt}{}

\titleformat{\chapter}[hang]{\LARGE\bfseries}{\thechapter\hsp\textcolor{blue}{|}\hsp}{0pt}{\Huge\bfseries}


\begin{document}
	
	\begin{titlepage}
		\begin{center}
			
			% Upper part of the page. The '~' is needed because \\
			% only works if a paragraph has started.
			\includegraphics[width=0.5\textwidth]{./img/logo.png}~\\[1cm]
			%?[width=0.15\textwidth]
			
			\textsc{\LARGE Politechnika Śląska w Gliwicach}\\[1.5cm]
			
			\textsc{\Large Obliczenia Równoległe II}\\[0.5cm]
			
			% Title
			\HRule \\[0.4cm]
			{ \huge \bfseries Komunikacja między procesami oraz redukcja danych (MPI) \\[0.4cm] }
			
			\HRule \\[1.5cm]
			
			% Author and supervisor
			\textsc{\Large Autor:} \\
			Bartłomiej Buchała \\
			[1.0cm]
			Informatyka SSM, semestr II \\
			Rok akademicki 2016/2017 \\
			Grupa OS1
			
			\vfill
			
			% Bottom of the page
			{\large \today}
			
		\end{center}
\end{titlepage}
	
	\newpage

	\bibliographystyle{plain}
\bibliography{bibliografia}	

\begin{thebibliography}{2}
	%\bibitem{Alias} Opis
	
	\bibitem{Czech} Z. J. Czech, \textit{Wprowadzenie do obliczeń równoległych}, Wydawnictwo PWN, Warszawa 2013, wyd. 2
	\bibitem{MPI} Message Passing Interface Forum, \textit{MPI: A Message-Passing Interface Standard, Version 3.0}, High Performance Computing Center Stuttgart (HLRS), Stuttgart 2012
	\bibitem{Pacheco} P. Pacheco, \textit{An Introduction to Parallel Programming}, Morgan Kaufmann, San Francisco 2001
	\bibitem{Parallel} M. J. Quinn, \textit{Parallel Programming in C with MPI and OpenMP}, McGraw-Hill, Nowy Jork 2003
	
\end{thebibliography}
\end{document}